\documentclass[11pt,a4paper]{article}

\usepackage[utf8]{inputenc}
\usepackage[
  left=1.05in,
  right=1.05in,
  top=1in,
  bottom=1in,
  headheight=14pt
]{geometry}
\usepackage{amsmath,amssymb,amsthm,mathtools}
\usepackage{bm}
\usepackage{enumitem}
\usepackage{booktabs,array}
\usepackage{microtype}
\usepackage{xcolor}
\usepackage[breakable,skins]{tcolorbox}
\usepackage[hidelinks]{hyperref}
\usepackage{fancyhdr}
\usepackage{etoolbox}
% Start a new page unless at least #1 of vertical space remains.
\makeatletter
\newcommand{\needspace}[1]{\par\penalty-100\begingroup
  \setlength{\dimen@}{#1}%
  \vskip\z@\@plus\dimen@\penalty-100\vskip\z@\@plus-\dimen@
  \vskip\dimen@\penalty9999\vskip-\dimen@\vskip\z@skip
  \endgroup}
\makeatother
\pretocmd{\section}{\needspace{12\baselineskip}}{}{}
\pretocmd{\subsection}{\needspace{8\baselineskip}}{}{}

\linespread{1.25}
\setlength{\parindent}{0pt}
\setlength{\parskip}{1em plus 0.25em minus 0.1em}
\setlength{\abovedisplayskip}{1.3em plus 0.3em minus 0.2em}
\setlength{\belowdisplayskip}{1.3em plus 0.3em minus 0.2em}
\setlength{\abovedisplayshortskip}{1em plus 0.2em}
\setlength{\belowdisplayshortskip}{1.1em plus 0.2em}
\setlength{\jot}{10pt}
\setlength{\emergencystretch}{2em}
\raggedbottom

\numberwithin{equation}{section}

\newtheoremstyle{spacedplain}
  {18pt}{18pt}{\itshape}{}{\bfseries}{.}{0.5em}{}
\newtheoremstyle{spaceddefinition}
  {18pt}{18pt}{\normalfont}{}{\bfseries}{.}{0.5em}{}
\newtheoremstyle{spacedremark}
  {18pt}{18pt}{\normalfont}{}{\itshape}{.}{0.5em}{}

\theoremstyle{spacedplain}
\newtheorem{theorem}{Theorem}[section]
\newtheorem{proposition}[theorem]{Proposition}

\theoremstyle{spaceddefinition}
\newtheorem{definition}[theorem]{Definition}
\newtheorem{exercise}[theorem]{Exercise}
\newtheorem{example}[theorem]{Example}

\theoremstyle{spacedremark}
\newtheorem{remark}[theorem]{Remark}

\setlist[itemize]{leftmargin=*,itemsep=7pt,topsep=8pt,parsep=3pt}
\setlist[enumerate]{leftmargin=*,itemsep=7pt,topsep=8pt,parsep=3pt}

% Plain black frame, used only for the most important results.
\newtcolorbox{keybox}{
  enhanced,
  colback=white,colframe=black,
  boxrule=0.6pt,sharp corners,
  left=10pt,right=10pt,top=4pt,bottom=6pt,
  before skip=16pt,after skip=16pt
}

\pagestyle{fancy}
\fancyhf{}
\fancyhead[L]{MTL712}
\fancyhead[R]{Lectures 35--40 and Hyperbolic Equations}
\fancyfoot[C]{\thepage}
\renewcommand{\headrulewidth}{0.4pt}
\renewcommand{\footrulewidth}{0pt}

% Run-in headings inside proofs and sections, set on their own line.
\newcommand{\step}[1]{\par\needspace{6\baselineskip}\addvspace{0.6em}\noindent\textbf{#1}\par\nobreak}
% Marks material that was not worked out in the lecture itself.
\newcommand{\note}[1]{\par\needspace{5\baselineskip}\addvspace{0.4em}\noindent\textit{(#1)}\par\nobreak}
\newcommand{\lecinfo}[1]{\noindent\textit{#1}\par}

\newcommand{\pd}[2]{\frac{\partial #1}{\partial #2}}
\newcommand{\pdd}[2]{\frac{\partial^2 #1}{\partial #2^2}}

\title{\textbf{MTL712: Computational Methods for Differential\\Equations}\\[6pt]
\large Lectures 35--40: Matrix Stability, Elliptic Equations,\\
ADI Iterations and Hyperbolic Equations}
\author{\normalsize Typeset lecture notes}
\date{\normalsize November 2026}

\begin{document}

\maketitle
\thispagestyle{plain}
{\linespread{1.0}\selectfont\setlength{\parskip}{0pt}\tableofcontents}

%=====================================================================
\section{Recap of Lecture 35}
%=====================================================================

\lecinfo{Lecture 35 is typeset in full in the previous notes; only the
results needed below are collected here.}

\step{Peaceman--Rachford ADI for \(u_t=u_{x_1x_1}+u_{x_2x_2}\)}
With \(r=k/h^2\), the two-dimensional Crank--Nicolson scheme
\[
  \bigl(1-\tfrac12r\delta_{x_1}^2\bigr)
  \bigl(1-\tfrac12r\delta_{x_2}^2\bigr)U_{l,m}^{n+1}
  =\bigl(1+\tfrac12r\delta_{x_1}^2\bigr)
   \bigl(1+\tfrac12r\delta_{x_2}^2\bigr)U_{l,m}^n
\]
is split into two steps, each implicit in one direction only:
\begin{align*}
  \bigl(1-\tfrac12r\delta_{x_1}^2\bigr)U_{l,m}^{n+\frac12}
  &=\bigl(1+\tfrac12r\delta_{x_2}^2\bigr)U_{l,m}^n,\\
  \bigl(1-\tfrac12r\delta_{x_2}^2\bigr)U_{l,m}^{n+1}
  &=\bigl(1+\tfrac12r\delta_{x_1}^2\bigr)U_{l,m}^{n+\frac12}.
\end{align*}

\step{Richardson's three-level scheme}
For \(u_t=u_{xx}\), central differences in time and space give
\[
  U_m^{n+1}=2r\,\delta_x^2U_m^n+U_m^{n-1}.
\]
Writing it as a two-level system for \(W=[U,\;V]^T\) with
\(V^n=U^{n-1}\), the amplification matrix is
\(\begin{bmatrix}-8r\sin^2\frac{\beta h}2&1\\ 1&0\end{bmatrix}\); one
of its eigenvalues has modulus greater than \(1\) for every \(r>0\), so
the scheme is \emph{unconditionally unstable}.

\step{Du Fort--Frankel scheme}
Replacing \(U_m^n\) in Richardson's scheme by
\(\frac12(U_m^{n+1}+U_m^{n-1})\) gives
\begin{equation}\label{eq:dff}
  U_m^{n+1}
  =\frac{2r}{1+2r}\bigl(U_{m+1}^n+U_{m-1}^n\bigr)
   +\frac{1-2r}{1+2r}\,U_m^{n-1}.
\end{equation}
Its stability is proved in Section~\ref{sec:dff}.

%=====================================================================
\section{Eigenvalues of Block Matrices}
%=====================================================================

\lecinfo{Lecture 36, 2 November 2026.}

\begin{keybox}
\begin{theorem}\label{thm:block}
Let
\[
  A=\begin{bmatrix}
    A_{11}&A_{12}&\cdots&A_{1m}\\
    A_{21}&A_{22}&\cdots&A_{2m}\\
    \vdots&&\ddots&\vdots\\
    A_{m1}&A_{m2}&\cdots&A_{mm}
  \end{bmatrix},
\]
where each \(A_{ij}\) is an \(n\times n\) matrix, and suppose all the
matrices \(A_{ij}\) have a common set of \(n\) linearly independent
eigenvectors \(v_1,\dots,v_n\).  Then the eigenvalues of \(A\) are
given by the eigenvalues of the \(m\times m\) matrices
\[
  \Lambda^{(k)}=
  \begin{bmatrix}
    \lambda_{11}^{(k)}&\lambda_{12}^{(k)}&\cdots&\lambda_{1m}^{(k)}\\
    \vdots&&\ddots&\vdots\\
    \lambda_{m1}^{(k)}&\lambda_{m2}^{(k)}&\cdots&\lambda_{mm}^{(k)}
  \end{bmatrix},
  \qquad k=1,2,\dots,n,
\]
where \(\lambda_{ij}^{(k)}\) is the eigenvalue of \(A_{ij}\)
corresponding to the \(k\)th common eigenvector \(v_k\).
\end{theorem}
\end{keybox}

\begin{proof}
Let \(v_k\) be an eigenvector common to all the submatrices
\(A_{ij}\), with corresponding eigenvalues
\(\lambda_{ij}^{(k)}\).  For simplicity take \(m=2\), i.e.\
\(i,j=1,2\), and write \(v\) for \(v_k\) and \(\lambda_{ij}\) for
\(\lambda_{ij}^{(k)}\).  Then
\[
  A_{11}v=\lambda_{11}v,\quad A_{12}v=\lambda_{12}v,\quad
  A_{21}v=\lambda_{21}v,\quad A_{22}v=\lambda_{22}v .
\]
Multiply these equations by non-zero constants \(\alpha_1,\alpha_2\)
(the first column by \(\alpha_1\), the second by \(\alpha_2\)) and
write them as
\begin{equation}\label{eq:bl1}
  \begin{bmatrix}A_{11}&A_{12}\\ A_{21}&A_{22}\end{bmatrix}
  \begin{bmatrix}\alpha_1v\\ \alpha_2v\end{bmatrix}
  =\begin{bmatrix}(\lambda_{11}\alpha_1+\lambda_{12}\alpha_2)v\\
     (\lambda_{21}\alpha_1+\lambda_{22}\alpha_2)v\end{bmatrix}.
\end{equation}
Assume that \(A\) has an eigenvalue \(\mu\) corresponding to the
eigenvector \(\begin{bmatrix}\alpha_1v\\ \alpha_2v\end{bmatrix}\), so
that
\begin{equation}\label{eq:bl2}
  \begin{bmatrix}A_{11}&A_{12}\\ A_{21}&A_{22}\end{bmatrix}
  \begin{bmatrix}\alpha_1v\\ \alpha_2v\end{bmatrix}
  =\mu\begin{bmatrix}\alpha_1v\\ \alpha_2v\end{bmatrix}.
\end{equation}
Comparing the right-hand sides of \eqref{eq:bl1} and \eqref{eq:bl2},
\begin{align*}
  \lambda_{11}\alpha_1+\lambda_{12}\alpha_2&=\mu\alpha_1,\\
  \lambda_{21}\alpha_1+\lambda_{22}\alpha_2&=\mu\alpha_2,
\end{align*}
that is,
\begin{align*}
  (\lambda_{11}-\mu)\alpha_1+\lambda_{12}\alpha_2&=0,\\
  \lambda_{21}\alpha_1+(\lambda_{22}-\mu)\alpha_2&=0.
\end{align*}
These two equations have a non-trivial solution for
\(\alpha_1,\alpha_2\) if and only if
\[
  \det\begin{bmatrix}\lambda_{11}-\mu&\lambda_{12}\\
    \lambda_{21}&\lambda_{22}-\mu\end{bmatrix}=0,
\]
i.e.\ if and only if \(\mu\) is an eigenvalue of
\(\begin{bmatrix}\lambda_{11}&\lambda_{12}\\
\lambda_{21}&\lambda_{22}\end{bmatrix}\).
\end{proof}

\begin{remark}
Each \(\Lambda^{(k)}\) has \(m\) eigenvalues, and \(k\) runs over
\(n\) eigenvectors, so this gives all \(mn\) eigenvalues of \(A\).
The argument for general \(m\) is identical, with \(m\) constants
\(\alpha_1,\dots,\alpha_m\).
\end{remark}

%=====================================================================
\section{Stability of the Du Fort--Frankel Scheme by the Matrix Method}
\label{sec:dff}
%=====================================================================

\lecinfo{Lecture 36, continued.}

Consider
\[
  \pd ut=\pdd ux,\qquad 0<x<1,\quad t>0,
\]
for a fixed mesh length \(h\) with \(Nh=1\), and \(u\) known on the
boundary \(x=0\), \(x=1\) for \(t>0\).  The Du Fort--Frankel scheme
\eqref{eq:dff} with \(r=k/h^2\) is
\[
  U_m^{n+1}=\frac{2r}{1+2r}\bigl(U_{m+1}^n+U_{m-1}^n\bigr)
  +\frac{1-2r}{1+2r}\,U_m^{n-1},
  \qquad m=1,\dots,N-1.
\]

\step{Matrix form}
For known boundary values these equations are
\begin{equation}\label{eq:dffmat}
  U^{n+1}=\frac{2r}{1+2r}\,AU^n+\frac{1-2r}{1+2r}\,U^{n-1}+f_n,
\end{equation}
where \(U^n=[U_1^n,\dots,U_{N-1}^n]^T\),
\[
  A=\begin{bmatrix}
    0&1&&&\\
    1&0&1&&\\
    &\ddots&\ddots&\ddots&\\
    &&1&0&1\\
    &&&1&0
  \end{bmatrix}_{(N-1)\times(N-1)},
  \qquad
  f_n=\frac{2r}{1+2r}
  \begin{bmatrix}U_0^n\\ 0\\ \vdots\\ 0\\ U_N^n\end{bmatrix}.
\]

\step{Reduction to a two-level scheme}
Define \(V^n:=U^{n-1}\).  Then \eqref{eq:dffmat} becomes
\[
  U^{n+1}=\frac{2r}{1+2r}\,AU^n+\frac{1-2r}{1+2r}\,V^n+f_n,
  \qquad
  V^{n+1}=U^n .
\]
Introduce \(W=\begin{bmatrix}U\\ V\end{bmatrix}\).  Then
\[
  W^{n+1}=BW^n+g^n,
  \qquad
  B=\begin{bmatrix}
    \dfrac{2r}{1+2r}A&\dfrac{1-2r}{1+2r}I\\[8pt]
    I&0
  \end{bmatrix},
  \qquad
  g^n=\begin{bmatrix}f_n\\ 0\end{bmatrix},
\]
where every block of \(B\) is \((N-1)\times(N-1)\).  This technique
has reduced a three-level difference scheme to a two-level one.

\begin{keybox}
The equations are stable if each eigenvalue of \(B\) has modulus
\(\le1\).
\end{keybox}

\step{Eigenvalues of \(B\)}
The eigenvalues of \(A\) are the \(N-1\) distinct numbers
\[
  \lambda_k=2\cos\frac{k\pi}N,\qquad k=1,2,\dots,N-1,
\]
so \(A\) possesses a set of \(N-1\) linearly independent eigenvectors
\(v_k\).  The same \(v_k\) are eigenvectors of \(I\) (eigenvalue
\(1\)) and of \(0\) (eigenvalue \(0\)).  Hence, by
Theorem~\ref{thm:block}, the eigenvalues \(\lambda\) of \(B\) are the
eigenvalues of
\[
  \begin{bmatrix}
    \dfrac{2r}{1+2r}\lambda_k&\dfrac{1-2r}{1+2r}\\[8pt]
    1&0
  \end{bmatrix},
  \qquad k=1,2,\dots,N-1.
\]
So
\[
  \det\begin{bmatrix}
    \dfrac{2r}{1+2r}\lambda_k-\lambda&\dfrac{1-2r}{1+2r}\\[8pt]
    1&-\lambda
  \end{bmatrix}=0
  \quad\Longrightarrow\quad
  \lambda^2-\frac{2r\lambda_k}{1+2r}\,\lambda-\frac{1-2r}{1+2r}=0 .
\]
Therefore
\[
  \lambda
  =\frac12\left[\frac{2r\lambda_k}{1+2r}
   \pm\left(\frac{4r^2\lambda_k^2}{(1+2r)^2}
   +\frac{4(1-2r)}{1+2r}\right)^{1/2}\right]
  =\frac{r\lambda_k\pm\bigl(r^2\lambda_k^2+1-4r^2\bigr)^{1/2}}{1+2r},
\]
that is,
\begin{equation}\label{eq:dfflam}
  \lambda=\frac{r\lambda_k\pm\bigl(1+(\lambda_k^2-4)r^2\bigr)^{1/2}}{1+2r}.
\end{equation}
Substituting \(\lambda_k=2\cos\frac{k\pi}N\) and using
\(4\cos^2\theta-4=-4\sin^2\theta\),
\[
  \lambda=\frac{2r\cos\frac{k\pi}N
  \pm\bigl(1-4r^2\sin^2\frac{k\pi}N\bigr)^{1/2}}{1+2r},
  \qquad k=1,2,\dots,N-1.
\]

\step{Case I: \(0\le1-4r^2\sin^2\frac{k\pi}N<1\)}
Both roots are real, and since \(|\cos\frac{k\pi}N|<1\),
\[
  |\lambda|<\frac{2r+1}{2r+1}=1 .
\]

\step{Case II: \(1-4r^2\sin^2\frac{k\pi}N<0\)}
The roots are complex conjugates, and
\[
  |\lambda|
  =\frac{\bigl(4r^2\cos^2\frac{k\pi}N+4r^2\sin^2\frac{k\pi}N-1\bigr)^{1/2}}{1+2r}
  =\frac{(4r^2-1)^{1/2}}{1+2r}
  =\left(\frac{4r^2-1}{4r^2+4r+1}\right)^{1/2}<1 .
\]

\begin{keybox}
\begin{theorem}
For a fixed mesh size, the Du Fort--Frankel scheme is unconditionally
stable for \(r>0\).
\end{theorem}
\end{keybox}

%=====================================================================
\section{Elliptic Equations}
%=====================================================================

\lecinfo{Lecture 37, 3 November 2026.}

Consider the second-order equation
\begin{equation}\label{eq:gen2}
  a(x_1,x_2)\pdd u{x_1}+2b(x_1,x_2)\frac{\partial^2u}{\partial x_1\partial x_2}
  +c(x_1,x_2)\pdd u{x_2}
  =d\Bigl(x_1,x_2,u,\pd u{x_1},\pd u{x_2}\Bigr),
\end{equation}
for \((x_1,x_2)\in R\), a bounded region in the \((x_1,x_2)\) plane
with boundary \(\partial R\).  It is \emph{elliptic} if
\[
  b^2-ac<0 .
\]

\subsection{Boundary value problems}

There are three distinct problems involving \eqref{eq:gen2},
depending on the boundary condition.

\begin{enumerate}[label=(\roman*)]
  \item \textbf{First BVP, or Dirichlet problem.}  The solution of
    \eqref{eq:gen2} takes prescribed values \(u(x_1,x_2)\) on the
    boundary \(\partial R\).
  \item \textbf{Second BVP, or Neumann problem.}
    \[
      \pd un=f(x_1,x_2)\quad\text{on }\partial R,
    \]
    where \(\partial/\partial n\) denotes differentiation along the
    normal to \(\partial R\) directed away from the interior of \(R\).
  \item \textbf{Third BVP, or Robin problem.}
    \[
      \alpha(x_1,x_2)\,u+\beta(x_1,x_2)\pd un=\gamma(x_1,x_2)
      \quad\text{on }\partial R,
    \]
    where \(\alpha(x_1,x_2)>0\) and \(\beta(x_1,x_2)>0\) for
    \((x_1,x_2)\in\partial R\).
\end{enumerate}

With \(a=1\), \(b=0\), \(c=1\) in \eqref{eq:gen2} we obtain
\[
  \pdd u{x_1}+\pdd u{x_2}=0\quad\text{(Laplace's equation)},
  \qquad
  \pdd u{x_1}+\pdd u{x_2}=f(x_1,x_2)\quad\text{(Poisson's equation)}.
\]

\subsection{The five-point formula}

Consider Laplace's equation
\[
  \pdd u{x_1}+\pdd u{x_2}=0
  \quad\text{on the unit square }
  R=[0,1]\times[0,1],
  \qquad u=g\ \text{on }\partial R .
\]
Use a square mesh \(x_{1,l}=lh\), \(x_{2,m}=mh\) with \(Nh=1\), so
\(x_{1,l+1}-x_{1,l}=x_{2,m+1}-x_{2,m}=h\).  There are \((N-1)^2\)
interior grid points.

\step{Taylor series}
\begin{align*}
  u_{l\pm1,m}&=u_{l,m}\pm h\pd{u_{l,m}}{x_1}
    +\frac{h^2}{2!}\pdd{u_{l,m}}{x_1}
    \pm\frac{h^3}{3!}\frac{\partial^3u_{l,m}}{\partial x_1^3}
    +\frac{h^4}{4!}\frac{\partial^4u_{l,m}}{\partial x_1^4}\pm\cdots,\\
  u_{l,m\pm1}&=u_{l,m}\pm h\pd{u_{l,m}}{x_2}
    +\frac{h^2}{2!}\pdd{u_{l,m}}{x_2}
    \pm\frac{h^3}{3!}\frac{\partial^3u_{l,m}}{\partial x_2^3}
    +\frac{h^4}{4!}\frac{\partial^4u_{l,m}}{\partial x_2^4}\pm\cdots.
\end{align*}
Adding the four expansions, the odd terms cancel:
\[
  u_{l+1,m}+u_{l-1,m}+u_{l,m+1}+u_{l,m-1}-4u_{l,m}
  =h^2\Bigl(\pdd{}{x_1}+\pdd{}{x_2}\Bigr)u_{l,m}
  +\frac{h^4}{12}\Bigl(\frac{\partial^4}{\partial x_1^4}
  +\frac{\partial^4}{\partial x_2^4}\Bigr)u_{l,m}+\cdots.
\]
So we can approximate
\[
  \Bigl(\pdd{}{x_1}+\pdd{}{x_2}\Bigr)u_{l,m}
  \approx\frac{u_{l+1,m}+u_{l-1,m}-4u_{l,m}+u_{l,m+1}+u_{l,m-1}}{h^2}.
\]
This gives a five-point analogue of the Laplacian, with local
truncation error
\[
  \frac{h^2}{12}\Bigl(\frac{\partial^4}{\partial x_1^4}
  +\frac{\partial^4}{\partial x_2^4}\Bigr)u_{l,m}+\cdots
\]
(the principal error term).

\begin{keybox}
\begin{definition}[Five-point formula]
The five-point finite difference replacement of Laplace's equation is
\[
  U_{l-1,m}+U_{l+1,m}-4U_{l,m}+U_{l,m+1}+U_{l,m-1}=0,
\]
a second-order approximation of the Laplacian, with
\(\mathrm{LTE}=O(h^2)\).
\end{definition}
\end{keybox}

\subsection{Matrix form}

Writing the five-point difference equation at all \((N-1)^2\) interior
grid points leads to a system
\[
  AU=b,
\]
where \(A\) is a block tridiagonal matrix of order \((N-1)^2\):
\[
  A=\begin{bmatrix}
    B&I&&\\
    I&B&I&\\
    &\ddots&\ddots&\ddots\\
    &&I&B
  \end{bmatrix},
  \qquad
  B=\begin{bmatrix}
    -4&1&&\\
    1&-4&1&\\
    &\ddots&\ddots&\ddots\\
    &&1&-4
  \end{bmatrix}_{(N-1)\times(N-1)},
\]
\(I\) is the identity matrix of order \(N-1\), and the unknown vector
is the \((N-1)^2\times1\) vector
\[
  U=[U_{1,1},U_{1,2},\dots,U_{1,N-1},U_{2,1},\dots,U_{2,N-1},
  \dots,U_{N-1,1},\dots,U_{N-1,N-1}]^T .
\]
The \((N-1)^2\times1\) vector \(b\) is known; it arises from the known
boundary values.

\begin{remark}
If we instead order \(U\) as
\([U_{1,1},U_{2,1},U_{3,1},\dots,U_{N-1,1},\dots,U_{N-1,N-1}]^T\), the
matrix has the same structure, since the five-point formula treats
\(x_1\) and \(x_2\) symmetrically.  Multiplying the equations by
\(-1\), \(A\) becomes \(\operatorname{tridiag}(-I,B,-I)\) with
\(B=\operatorname{tridiag}(-1,4,-1)\).
\end{remark}

\begin{exercise}\label{ex:eigA}
Show that the eigenvalues of \(A\) are
\[
  \lambda_{ij}=-4\Bigl(\sin^2\frac{i\pi}{2N}+\sin^2\frac{j\pi}{2N}\Bigr),
  \qquad 1\le i,j\le N-1,
\]
and hence \(|A|\ne0\).
\end{exercise}

\note{This was an exercise; solution not done in lecture.}

\begin{proof}[Solution]
The tridiagonal matrix \(\operatorname{tridiag}(1,0,1)\) of order
\(N-1\) has eigenvalues \(2\cos\frac{j\pi}N\) with eigenvectors
\(v_j=\bigl[\sin\frac{j\pi}N,\sin\frac{2j\pi}N,\dots,
\sin\frac{(N-1)j\pi}N\bigr]^T\).  Since \(B=-4I+\operatorname{tridiag}(1,0,1)\),
the blocks \(B\) and \(I\) share the eigenvectors \(v_i\), with
eigenvalues
\[
  \beta_i=-4+2\cos\frac{i\pi}N\quad\text{for }B,
  \qquad 1\quad\text{for }I.
\]
By Theorem~\ref{thm:block}, the eigenvalues of \(A\) are the
eigenvalues of the \((N-1)\times(N-1)\) matrices
\[
  \operatorname{tridiag}(1,\beta_i,1)=\beta_iI+\operatorname{tridiag}(1,0,1),
\]
namely \(\beta_i+2\cos\frac{j\pi}N\).  Thus
\[
  \lambda_{ij}=-4+2\cos\frac{i\pi}N+2\cos\frac{j\pi}N .
\]
Using \(2\cos\theta-2=-4\sin^2\frac\theta2\),
\[
  \lambda_{ij}=-4\sin^2\frac{i\pi}{2N}-4\sin^2\frac{j\pi}{2N}.
\]
For \(1\le i,j\le N-1\) every \(\lambda_{ij}<0\), so no eigenvalue is
zero and \(|A|\ne0\): the system \(AU=b\) has a unique solution.
\end{proof}

\subsection{Rectangular mesh}

\begin{exercise}\label{ex:rect}
For \(\pdd u{x_1}+\pdd u{x_2}=0\) on a rectangle with step lengths
\(h_1\) and \(h_2\) in the \(x_1\) and \(x_2\) directions
respectively, derive the five-point difference replacement at the grid
point \((lh_1,mh_2)\):
\[
  \theta^{-1}\bigl(U_{l+1,m}+U_{l-1,m}\bigr)
  +\theta\bigl(U_{l,m+1}+U_{l,m-1}\bigr)
  -2\bigl(\theta+\theta^{-1}\bigr)U_{l,m}=0,
  \qquad
  \theta=\frac{h_1}{h_2}\ \text{(the mesh ratio)}.
\]
Show that the local truncation error of this formula is
\[
  \mathrm{LTE}=\frac1{12}\Bigl(h_1^2\frac{\partial^4u}{\partial x_1^4}
  +h_2^2\frac{\partial^4u}{\partial x_2^4}\Bigr)+\cdots.
\]
\end{exercise}

\note{This was an exercise; solution not done in lecture.}

\begin{proof}[Solution]
\step{Derivation}
Taylor expansion in \(x_1\) with step \(h_1\) gives
\[
  u_{l+1,m}+u_{l-1,m}-2u_{l,m}
  =h_1^2\pdd u{x_1}+\frac{h_1^4}{12}\frac{\partial^4u}{\partial x_1^4}
  +O(h_1^6),
\]
and similarly in \(x_2\) with step \(h_2\):
\[
  u_{l,m+1}+u_{l,m-1}-2u_{l,m}
  =h_2^2\pdd u{x_2}+\frac{h_2^4}{12}\frac{\partial^4u}{\partial x_2^4}
  +O(h_2^6).
\]
Replacing each second derivative by its central difference, Laplace's
equation becomes
\[
  \frac{U_{l+1,m}-2U_{l,m}+U_{l-1,m}}{h_1^2}
  +\frac{U_{l,m+1}-2U_{l,m}+U_{l,m-1}}{h_2^2}=0 .
\]
Multiplying by \(h_1h_2\) and using \(h_2/h_1=\theta^{-1}\),
\(h_1/h_2=\theta\),
\[
  \theta^{-1}\bigl(U_{l+1,m}-2U_{l,m}+U_{l-1,m}\bigr)
  +\theta\bigl(U_{l,m+1}-2U_{l,m}+U_{l,m-1}\bigr)=0,
\]
which is the required formula after collecting the \(U_{l,m}\) terms.

\step{Local truncation error}
Substitute the exact solution \(u\) into the left-hand side and divide
by \(h_1h_2\) (so that the operator approximates
\(\pdd u{x_1}+\pdd u{x_2}\)):
\begin{align*}
  &\frac1{h_1h_2}\Bigl[\theta^{-1}\bigl(u_{l+1,m}+u_{l-1,m}-2u_{l,m}\bigr)
   +\theta\bigl(u_{l,m+1}+u_{l,m-1}-2u_{l,m}\bigr)\Bigr]\\
  &\qquad=\frac1{h_1^2}\Bigl(h_1^2\pdd u{x_1}
     +\frac{h_1^4}{12}\frac{\partial^4u}{\partial x_1^4}\Bigr)
   +\frac1{h_2^2}\Bigl(h_2^2\pdd u{x_2}
     +\frac{h_2^4}{12}\frac{\partial^4u}{\partial x_2^4}\Bigr)+\cdots\\
  &\qquad=\Bigl(\pdd u{x_1}+\pdd u{x_2}\Bigr)
   +\frac1{12}\Bigl(h_1^2\frac{\partial^4u}{\partial x_1^4}
   +h_2^2\frac{\partial^4u}{\partial x_2^4}\Bigr)+\cdots.
\end{align*}
Since \(u\) satisfies Laplace's equation, the first bracket vanishes,
leaving
\[
  \mathrm{LTE}=\frac1{12}\Bigl(h_1^2\frac{\partial^4u}{\partial x_1^4}
  +h_2^2\frac{\partial^4u}{\partial x_2^4}\Bigr)+\cdots
  =O(h_1^2+h_2^2).
\]
For \(h_1=h_2=h\) (\(\theta=1\)) this reduces to the five-point formula
and its error term above.
\end{proof}

%=====================================================================
\section{ADI Iteration for Laplace's Equation: Dirichlet Problem}
%=====================================================================

\lecinfo{Lectures 38--39.}

\subsection{The Peaceman--Rachford method as an iteration}

Consider
\[
  \pd ut=\pdd u{x_1}+\pdd u{x_2}
\]
with appropriate initial and boundary conditions.  The
Peaceman--Rachford ADI method is
\begin{equation}\label{eq:pr}
\begin{aligned}
  \bigl(1-\tfrac12\lambda\delta_{x_1}^2\bigr)U_{l,m}^{n+1*}
  &=\bigl(1+\tfrac12\lambda\delta_{x_2}^2\bigr)U_{l,m}^n,\\
  \bigl(1-\tfrac12\lambda\delta_{x_2}^2\bigr)U_{l,m}^{n+1}
  &=\bigl(1+\tfrac12\lambda\delta_{x_1}^2\bigr)U_{l,m}^{n+1*},
\end{aligned}
\end{equation}
where \(U_{l,m}^{n+1*}=U_{l,m}^{n+1/2}\).  Eliminating
\(U_{l,m}^{n+1*}\) from \eqref{eq:pr} we get
\begin{equation}\label{eq:prc}
  \bigl(1-\tfrac12\lambda\delta_{x_1}^2\bigr)
  \bigl(1-\tfrac12\lambda\delta_{x_2}^2\bigr)U_{l,m}^{n+1}
  =\bigl(1+\tfrac12\lambda\delta_{x_1}^2\bigr)
  \bigl(1+\tfrac12\lambda\delta_{x_2}^2\bigr)U_{l,m}^n .
\end{equation}

Treat \eqref{eq:prc} as an iteration scheme, which is assumed to be
convergent, i.e.\ \(U_{l,m}^{n+1}=U_{l,m}^n=U_{l,m}\) for \(n\)
sufficiently large.  Then
\[
  \bigl(1-\tfrac12\lambda\delta_{x_1}^2\bigr)
  \bigl(1-\tfrac12\lambda\delta_{x_2}^2\bigr)U_{l,m}
  =\bigl(1+\tfrac12\lambda\delta_{x_1}^2\bigr)
  \bigl(1+\tfrac12\lambda\delta_{x_2}^2\bigr)U_{l,m},
\]
that is,
\[
  \Bigl(1-\tfrac\lambda2\bigl(\delta_{x_1}^2+\delta_{x_2}^2\bigr)
  +\tfrac{\lambda^2}4\delta_{x_1}^2\delta_{x_2}^2\Bigr)U_{l,m}
  =\Bigl(1+\tfrac\lambda2\bigl(\delta_{x_1}^2+\delta_{x_2}^2\bigr)
  +\tfrac{\lambda^2}4\delta_{x_1}^2\delta_{x_2}^2\Bigr)U_{l,m}.
\]
Hence
\[
  \lambda\bigl(\delta_{x_1}^2+\delta_{x_2}^2\bigr)U_{l,m}=0
  \quad\Longrightarrow\quad
  \bigl(\delta_{x_1}^2+\delta_{x_2}^2\bigr)U_{l,m}=0,
\]
which is the five-point analogue of Laplace's equation.

\begin{keybox}
The Peaceman--Rachford ADI method for the heat conduction problem,
with boundary conditions independent of time, is an iteration method
for solving Laplace's equation subject to the Dirichlet condition.
Here \(\lambda\) is no longer a mesh parameter; it is an
\emph{iteration parameter}.
\end{keybox}

\subsection{Matrix form}

Let \(Nh=1\) and order the unknowns at the \((N-1)^2\) interior grid
points with the \(x_1\) index running fastest:
\[
  U=[U_{1,1},U_{2,1},\dots,U_{N-1,1},\;U_{1,2},\dots,U_{N-1,2},\;
  \dots,\;U_{1,N-1},\dots,U_{N-1,N-1}]^T,
\]
an \((N-1)^2\times1\) unknown vector.  Let \(H\) and \(V\) be the
matrices of \(-\delta_{x_1}^2\) and \(-\delta_{x_2}^2\) acting on
interior values:
\[
  H=\begin{bmatrix}C&&\\ &\ddots&\\ &&C\end{bmatrix},
  \qquad
  C=\begin{bmatrix}
    2&-1&&\\
    -1&2&-1&\\
    &\ddots&\ddots&\ddots\\
    &&-1&2
  \end{bmatrix}_{(N-1)\times(N-1)},
\]
so \(H\) is block diagonal of order \((N-1)^2\), and
\[
  V=\begin{bmatrix}
    2J&-J&&\\
    -J&2J&-J&\\
    &\ddots&\ddots&\ddots\\
    &&-J&2J
  \end{bmatrix},
\]
a block tridiagonal matrix, where \(J\) is the identity matrix of
order \(N-1\).  The vector
\(K=[K_{1,1},\dots,K_{N-1,1},\dots,K_{1,N-1},\dots,K_{N-1,N-1}]^T\) is
known; it arises from the boundary values.  The five-point equations at
all interior points read
\[
  (H+V)\,U=K .
\]

Dividing \eqref{eq:pr} by \(\lambda/2\), the totality of the equations
at all \((N-1)^2\) interior grid points leads to the matrix--vector
equations
\begin{keybox}
\begin{equation}\label{eq:adimat}
\begin{aligned}
  \Bigl(\frac2{\lambda_{n+1}}I+H\Bigr)U^{n+1*}
  &=\Bigl(\frac2{\lambda_{n+1}}I-V\Bigr)U^n+K,\\
  \Bigl(\frac2{\lambda_{n+1}}I+V\Bigr)U^{n+1}
  &=\Bigl(\frac2{\lambda_{n+1}}I-H\Bigr)U^{n+1*}+K,
\end{aligned}
\end{equation}
\end{keybox}
where the parameter \(\lambda_{n+1}\) can vary from iteration to
iteration.

\step{Limit of the iteration}
Eliminating \(U^{n+1*}\),
\[
  U^{n+1*}=\Bigl(\frac2{\lambda_{n+1}}I+H\Bigr)^{-1}
  \Bigl[\Bigl(\frac2{\lambda_{n+1}}I-V\Bigr)U^n+K\Bigr],
\]
\[
  \Bigl(\frac2{\lambda_{n+1}}I+V\Bigr)U^{n+1}
  =\Bigl(\frac2{\lambda_{n+1}}I-H\Bigr)
  \Bigl(\frac2{\lambda_{n+1}}I+H\Bigr)^{-1}
  \Bigl[\Bigl(\frac2{\lambda_{n+1}}I-V\Bigr)U^n+K\Bigr]+K .
\]
Multiplying by \(\bigl(\frac2{\lambda_{n+1}}I+H\bigr)\), which commutes
with \(\bigl(\frac2{\lambda_{n+1}}I-H\bigr)\),
\begin{align*}
  \Bigl(\frac2{\lambda_{n+1}}I+H\Bigr)\Bigl(\frac2{\lambda_{n+1}}I+V\Bigr)U^{n+1}
  &=\Bigl(\frac2{\lambda_{n+1}}I-H\Bigr)\Bigl(\frac2{\lambda_{n+1}}I-V\Bigr)U^n\\
  &\quad+\Bigl(\frac2{\lambda_{n+1}}I-H\Bigr)K+\Bigl(\frac2{\lambda_{n+1}}I+H\Bigr)K .
\end{align*}
Assuming \(U^{n+1}=U^n=U\) and expanding both products,
\[
  \Bigl[\frac4{\lambda_{n+1}^2}I+\frac2{\lambda_{n+1}}(H+V)+HV\Bigr]U
  =\Bigl[\frac4{\lambda_{n+1}^2}I-\frac2{\lambda_{n+1}}(H+V)+HV\Bigr]U
  +\frac4{\lambda_{n+1}}K,
\]
so
\[
  \frac4{\lambda_{n+1}}(H+V)\,U=\frac4{\lambda_{n+1}}K
  \quad\Longrightarrow\quad
  (H+V)\,U=K .
\]
The limit is the solution of the five-point equations.

\subsection{Convergence}

The iteration scheme \eqref{eq:adimat} can be written as
\[
  U^{n+1}=T_{n+1}U^n+g_{n+1},
\]
where
\begin{equation}\label{eq:T}
  T_{n+1}=\Bigl(\frac2{\lambda_{n+1}}I+V\Bigr)^{-1}
  \Bigl(\frac2{\lambda_{n+1}}I+H\Bigr)^{-1}
  \Bigl[\Bigl(\frac2{\lambda_{n+1}}I-H\Bigr)
  \Bigl(\frac2{\lambda_{n+1}}I-V\Bigr)\Bigr]
\end{equation}
and
\[
  g_{n+1}=\Bigl(\frac2{\lambda_{n+1}}I+V\Bigr)^{-1}
  \Bigl[\Bigl(\frac2{\lambda_{n+1}}I+H\Bigr)^{-1}
  \Bigl(\frac2{\lambda_{n+1}}I-H\Bigr)+I\Bigr]K .
\]

Since the matrices \(H\) and \(V\) commute, they possess a common set
of orthonormal eigenvectors \(\alpha_{p,q}\), \(1\le p,q\le N-1\),
with
\[
  H\alpha_{p,q}=4\sin^2\frac{p\pi}{2N}\,\alpha_{p,q},
  \qquad
  V\alpha_{p,q}=4\sin^2\frac{q\pi}{2N}\,\alpha_{p,q},
  \qquad 1\le p,q\le N-1 .
\]
(The eigenvalues of \(C\) are \(4\sin^2\frac{p\pi}{2N}\).)  From
\eqref{eq:T},
\begin{align*}
  T_{n+1}\alpha_{p,q}
  &=\frac{\Bigl(\frac2{\lambda_{n+1}}-4\sin^2\frac{p\pi}{2N}\Bigr)
          \Bigl(\frac2{\lambda_{n+1}}-4\sin^2\frac{q\pi}{2N}\Bigr)}
         {\Bigl(\frac2{\lambda_{n+1}}+4\sin^2\frac{p\pi}{2N}\Bigr)
          \Bigl(\frac2{\lambda_{n+1}}+4\sin^2\frac{q\pi}{2N}\Bigr)}
     \,\alpha_{p,q}\\[4pt]
  &=\frac{\bigl(1-2\lambda_{n+1}\sin^2\frac{p\pi}{2N}\bigr)
          \bigl(1-2\lambda_{n+1}\sin^2\frac{q\pi}{2N}\bigr)}
         {\bigl(1+2\lambda_{n+1}\sin^2\frac{p\pi}{2N}\bigr)
          \bigl(1+2\lambda_{n+1}\sin^2\frac{q\pi}{2N}\bigr)}
     \,\alpha_{p,q}.
\end{align*}
Let
\[
  a=2\lambda_{n+1}\sin^2\frac{p\pi}{2N},
  \qquad
  b=2\lambda_{n+1}\sin^2\frac{q\pi}{2N},
  \qquad a>0,\ b>0 .
\]
Then the eigenvalue is
\[
  \frac{(1-a)(1-b)}{(1+a)(1+b)}
  =\frac{1-a-b+ab}{1+a+b+ab}
  =1-\frac{2(a+b)}{(1+a)(1+b)} .
\]
Under the conditions on \(a,b\),
\[
  -1<1-\frac{2(a+b)}{(1+a)(1+b)}<1\qquad\forall\,a,b>0 .
\]
(The right inequality holds because \(a+b>0\); the left one because
\(a+b<1+a+b+ab\).)  Hence \(\rho(T_{n+1})<1\) for all \(a,b\), and the
iteration converges.

\begin{keybox}
\begin{theorem}
For every choice of the iteration parameters \(\lambda_{n+1}>0\), the
Peaceman--Rachford iteration \eqref{eq:adimat} converges to the
solution of the five-point equations \((H+V)U=K\).
\end{theorem}
\end{keybox}

\subsection{A variant: exercise}

\begin{exercise}
Do the convergence analysis for the scheme
\begin{align*}
  \Bigl(\frac1{\lambda_{n+1}}I+H\Bigr)U^{n+1*}
  &=\Bigl(\frac1{\lambda_{n+1}}I-V\Bigr)U^n+K,\\
  \Bigl(\frac1{\lambda_{n+1}}I+V\Bigr)U^{n+1}
  &=\Bigl(\frac1{\lambda_{n+1}}I-H\Bigr)U^{n+1*}+K,
\end{align*}
with \(H\) and \(V\) as above.
\end{exercise}

\note{This was an exercise; solution not done in lecture.}

\begin{proof}[Solution]
This is \eqref{eq:adimat} with \(2/\lambda_{n+1}\) replaced by
\(1/\lambda_{n+1}\).  Exactly as before, the fixed point satisfies
\((H+V)U=K\), and the iteration matrix
\[
  T_{n+1}=\Bigl(\frac1{\lambda_{n+1}}I+V\Bigr)^{-1}
  \Bigl(\frac1{\lambda_{n+1}}I+H\Bigr)^{-1}
  \Bigl(\frac1{\lambda_{n+1}}I-H\Bigr)\Bigl(\frac1{\lambda_{n+1}}I-V\Bigr)
\]
has the eigenvalues
\[
  \frac{(1-a)(1-b)}{(1+a)(1+b)},
  \qquad
  a=4\lambda_{n+1}\sin^2\frac{p\pi}{2N},\quad
  b=4\lambda_{n+1}\sin^2\frac{q\pi}{2N},
\]
with \(a,b>0\).  The same argument gives modulus \(<1\), so this
iteration also converges for every \(\lambda_{n+1}>0\).
\end{proof}

%=====================================================================
\section{The Douglas--Rachford ADI Method}
%=====================================================================

\lecinfo{Lecture 39.}

For
\[
  \pd ut=\pdd u{x_1}+\pdd u{x_2},
  \qquad (x_1,x_2,t)\in[0,1]^2\times[0,\infty),
\]
with initial and boundary conditions, consider
\begin{equation}\label{eq:dr}
  \bigl(1-\lambda\delta_{x_1}^2\bigr)\bigl(1-\lambda\delta_{x_2}^2\bigr)U_{l,m}^{n+1}
  =\bigl(1+\lambda^2\delta_{x_1}^2\delta_{x_2}^2\bigr)U_{l,m}^n .
\end{equation}

\step{As an iteration method}
Treat \eqref{eq:dr} as an iteration which is convergent:
\(U_{l,m}^{n+1}\approx U_{l,m}^n\approx U_{l,m}\).  Then
\[
  \bigl(1-\lambda\delta_{x_1}^2\bigr)\bigl(1-\lambda\delta_{x_2}^2\bigr)U_{l,m}
  =\bigl(1+\lambda^2\delta_{x_1}^2\delta_{x_2}^2\bigr)U_{l,m},
\]
\[
  \Bigl(1-\lambda\bigl(\delta_{x_1}^2+\delta_{x_2}^2\bigr)
  +\lambda^2\delta_{x_1}^2\delta_{x_2}^2\Bigr)U
  =\bigl(1+\lambda^2\delta_{x_1}^2\delta_{x_2}^2\bigr)U,
\]
so \(\lambda\bigl(\delta_{x_1}^2+\delta_{x_2}^2\bigr)U=0\), i.e.\
\(\bigl(\delta_{x_1}^2+\delta_{x_2}^2\bigr)U=0\): again the five-point
Laplace equation.

\begin{keybox}
\begin{definition}[Douglas--Rachford ADI, split form]
\begin{align*}
  \bigl(1-\lambda\delta_{x_1}^2\bigr)U_{l,m}^{n+1*}
  &=\bigl(1+\lambda\delta_{x_2}^2\bigr)U_{l,m}^n,\\
  \bigl(1-\lambda\delta_{x_2}^2\bigr)U_{l,m}^{n+1}
  &=U_{l,m}^{n+1*}-\lambda\delta_{x_2}^2U_{l,m}^n .
\end{align*}
\end{definition}
\end{keybox}

\step{Check}
From the second equation,
\(U^{n+1*}=\bigl(1-\lambda\delta_{x_2}^2\bigr)U^{n+1}+\lambda\delta_{x_2}^2U^n\).
Applying \(\bigl(1-\lambda\delta_{x_1}^2\bigr)\) and using the first
equation,
\[
  \bigl(1-\lambda\delta_{x_1}^2\bigr)\bigl(1-\lambda\delta_{x_2}^2\bigr)U^{n+1}
  =\bigl(1+\lambda\delta_{x_2}^2\bigr)U^n
   -\lambda\delta_{x_2}^2\bigl(1-\lambda\delta_{x_1}^2\bigr)U^n
  =\bigl(1+\lambda^2\delta_{x_1}^2\delta_{x_2}^2\bigr)U^n,
\]
which is \eqref{eq:dr}.

%=====================================================================
\section{ADI Iteration for the Neumann Problem}
%=====================================================================

\lecinfo{Lectures 39--40.}

\subsection{The problem and its discretization}

Consider
\[
  \pdd u{x_1}+\pdd u{x_2}=0,\qquad 0\le x_1,x_2\le1,
\]
\[
  \pd un=g(x_1,x_2),\qquad (x_1,x_2)\in\partial R,
\]
where \(n\) is the outwardly directed normal.  On the four sides,
\begin{align*}
  -\pd u{x_1}&=g(0,x_2)\ \text{along }x_1=0,
  &\pd u{x_1}&=g(1,x_2)\ \text{along }x_1=1,
  &&0\le x_2\le1,\\
  -\pd u{x_2}&=g(x_1,0)\ \text{along }x_2=0,
  &\pd u{x_2}&=g(x_1,1)\ \text{along }x_2=1,
  &&0\le x_1\le1,
\end{align*}
where \(\partial u/\partial x_1\) and \(\partial u/\partial x_2\) are
taken in the increasing directions of \(x_1\) and \(x_2\).  Take
\(Nh=1\).  Since \(u\) is not known on the boundary, \emph{all}
\((N+1)^2\) grid points \((l,m)\), \(0\le l,m\le N\), are unknowns.

\step{Fictitious points}
Apply the five-point formula at a boundary point, using fictitious
points outside the square.  At the corner \((0,0)\):
\[
  \frac{u_{1,0}+u_{-1,0}+u_{0,1}+u_{0,-1}-4u_{0,0}}{h^2}=0 .
\]
The boundary condition along \(x_1=0\), by a central difference, gives
\[
  -\frac{u_{1,0}-u_{-1,0}}{2h}=g_{0,0}
  \quad\Longrightarrow\quad
  u_{-1,0}=2hg_{0,0}+u_{1,0},
\]
and along \(x_2=0\),
\[
  -\frac{u_{0,1}-u_{0,-1}}{2h}=g_{0,0}
  \quad\Longrightarrow\quad
  u_{0,-1}=2hg_{0,0}+u_{0,1}.
\]
Substituting,
\[
  4u_{0,0}-2u_{1,0}-2u_{0,1}=2h\cdot2g_{0,0}.
\]
Similarly, at a non-corner point \((0,m)\) of the side \(x_1=0\),
\[
  4u_{0,m}-2u_{1,m}-u_{0,m+1}-u_{0,m-1}=2h\,g_{0,m},
\]
and the same holds on the other sides.

\step{The difference system}
The five-point analogue of the Laplacian,
\[
  \pdd u{x_1}+\pdd u{x_2}
  \approx\frac1{h^2}\bigl(u_{l+1,m}+u_{l-1,m}+u_{l,m+1}+u_{l,m-1}-4u_{l,m}\bigr),
\]
together with the boundary relations, at all \((N+1)^2\) grid points
gives rise to
\begin{keybox}
\[
  AU=2hG,
\]
\end{keybox}
where \(A\) is a block tridiagonal matrix of order \((N+1)^2\):
\[
  A=\begin{bmatrix}
    B&-2I&&&\\
    -I&B&-I&&\\
    &\ddots&\ddots&\ddots&\\
    &&-I&B&-I\\
    &&&-2I&B
  \end{bmatrix},
  \qquad
  B=\begin{bmatrix}
    4&-2&&&\\
    -1&4&-1&&\\
    &\ddots&\ddots&\ddots&\\
    &&-1&4&-1\\
    &&&-2&4
  \end{bmatrix}_{(N+1)\times(N+1)},
\]
with
\[
  U=[u_{0,0},u_{1,0},\dots,u_{N,0},\;u_{0,1},\dots,u_{N,1},\;\dots,\;
  u_{0,N},\dots,u_{N,N}]^T,
\]
and
\[
  G=[2g_{0,0},g_{1,0},\dots,g_{N-1,0},2g_{N,0},\;g_{0,1},0,\dots,0,g_{N,1},\;
  \dots,\;2g_{0,N},g_{1,N},\dots,g_{N-1,N},2g_{N,N}]^T .
\]
That is, the entry of \(G\) is \(2g\) at the four corners, \(g\) at
the other boundary points, and \(0\) at interior points.

\begin{remark}
\note{Added beyond the notes.}
Every row of \(A\) sums to zero, so \(A\mathbf 1=0\): \(A\) is singular.
This reflects the fact that the solution of the Neumann problem is
determined only up to an additive constant.  A solution exists only if
the data satisfy a compatibility condition, the discrete form of
\(\oint_{\partial R}g\,ds=0\).
\end{remark}

\subsection{The Peaceman--Rachford iteration}

Split \(A=X+Y\), where \(X\) and \(Y\) contain the \(x_1\) and \(x_2\)
differences respectively:
\[
  X=\begin{bmatrix}\tilde C&&\\ &\ddots&\\ &&\tilde C\end{bmatrix},
  \qquad
  Y=\begin{bmatrix}
    2I&-2I&&\\
    -I&2I&-I&\\
    &\ddots&\ddots&\ddots\\
    &&-2I&2I
  \end{bmatrix},
\]
where
\[
  \tilde C=\begin{bmatrix}
    2&-2&&\\
    -1&2&-1&\\
    &\ddots&\ddots&\ddots\\
    &&-2&2
  \end{bmatrix}_{(N+1)\times(N+1)}.
\]
The Peaceman--Rachford iteration is then
\begin{keybox}
\begin{equation}\label{eq:neuadi}
\begin{aligned}
  \Bigl(\frac2{\lambda_{n+1}}I+X\Bigr)U^{n+1*}
  &=\Bigl(\frac2{\lambda_{n+1}}I-Y\Bigr)U^n+2hG,\\
  \Bigl(\frac2{\lambda_{n+1}}I+Y\Bigr)U^{n+1}
  &=\Bigl(\frac2{\lambda_{n+1}}I-X\Bigr)U^{n+1*}+2hG .
\end{aligned}
\end{equation}
\end{keybox}

\step{Eliminating \(U^{n+1*}\)}
\[
  U^{n+1*}=\Bigl(\frac2{\lambda_{n+1}}I+X\Bigr)^{-1}
  \Bigl\{\Bigl(\frac2{\lambda_{n+1}}I-Y\Bigr)U^n+2hG\Bigr\},
\]
and, since \(\bigl(\frac2{\lambda_{n+1}}I-X\bigr)\) and
\(\bigl(\frac2{\lambda_{n+1}}I+X\bigr)^{-1}\) commute,
\begin{equation}\label{eq:neuS}
\begin{aligned}
  U^{n+1}
  &=\Bigl(\frac2{\lambda_{n+1}}I+Y\Bigr)^{-1}
    \Bigl(\frac2{\lambda_{n+1}}I+X\Bigr)^{-1}
    \Bigl(\frac2{\lambda_{n+1}}I-X\Bigr)
    \Bigl(\frac2{\lambda_{n+1}}I-Y\Bigr)U^n\\
  &\quad+\Bigl(\frac2{\lambda_{n+1}}I+Y\Bigr)^{-1}
    \Bigl[\Bigl(\frac2{\lambda_{n+1}}I+X\Bigr)^{-1}
    \Bigl(\frac2{\lambda_{n+1}}I-X\Bigr)+I\Bigr]2hG .
\end{aligned}
\end{equation}

\step{Limit of the iteration}
Put \(U^{n+1}=U^n=U\) (convergence).  Multiplying by
\(\bigl(\frac2{\lambda_{n+1}}I+X\bigr)\bigl(\frac2{\lambda_{n+1}}I+Y\bigr)\),
\begin{align*}
  \Bigl(\frac2{\lambda_{n+1}}I+X\Bigr)\Bigl(\frac2{\lambda_{n+1}}I+Y\Bigr)U
  &=\Bigl(\frac2{\lambda_{n+1}}I-X\Bigr)\Bigl(\frac2{\lambda_{n+1}}I-Y\Bigr)U\\
  &\quad+\Bigl[\Bigl(\frac2{\lambda_{n+1}}I-X\Bigr)+\Bigl(\frac2{\lambda_{n+1}}I+X\Bigr)\Bigr]2hG .
\end{align*}
Expanding,
\[
  \frac4{\lambda_{n+1}^2}U+\frac2{\lambda_{n+1}}(X+Y)U+XYU
  =\frac4{\lambda_{n+1}^2}U-\frac2{\lambda_{n+1}}(X+Y)U+XYU
  +\frac4{\lambda_{n+1}}(2hG),
\]
so
\[
  \frac4{\lambda_{n+1}}(X+Y)U=\frac4{\lambda_{n+1}}\,2hG
  \quad\Longrightarrow\quad
  (X+Y)U=2hG,\quad\text{i.e.}\quad AU=2hG .
\]

\subsection{Convergence}

Equation \eqref{eq:neuS} can be written as
\[
  U^{n+1}=S_{n+1}U^n+g_{n+1},
\]
where
\[
  S_{n+1}=\Bigl(\frac2{\lambda_{n+1}}I+Y\Bigr)^{-1}
  \Bigl(\frac2{\lambda_{n+1}}I+X\Bigr)^{-1}
  \Bigl(\frac2{\lambda_{n+1}}I-X\Bigr)
  \Bigl(\frac2{\lambda_{n+1}}I-Y\Bigr),
\]
and similarly for \(g_{n+1}\).

Since \(XY=YX\), these matrices possess a common set of linearly
independent eigenvectors \(\beta_{p,q}\), \(0\le p,q\le N\),
corresponding to
\[
  X\beta_{p,q}=4\sin^2\frac{p\pi}{2N}\,\beta_{p,q},
  \qquad
  Y\beta_{p,q}=4\sin^2\frac{q\pi}{2N}\,\beta_{p,q}.
\]
(The eigenvalues of \(\tilde C\) are \(4\sin^2\frac{p\pi}{2N}\),
\(p=0,1,\dots,N\), with eigenvectors having components
\(\cos\frac{lp\pi}N\), \(l=0,\dots,N\).)  Hence
\[
  S_{n+1}\beta_{p,q}
  =\frac{\Bigl(\frac2{\lambda_{n+1}}-4\sin^2\frac{p\pi}{2N}\Bigr)
         \Bigl(\frac2{\lambda_{n+1}}-4\sin^2\frac{q\pi}{2N}\Bigr)}
        {\Bigl(\frac2{\lambda_{n+1}}+4\sin^2\frac{p\pi}{2N}\Bigr)
         \Bigl(\frac2{\lambda_{n+1}}+4\sin^2\frac{q\pi}{2N}\Bigr)}\,\beta_{p,q}
  =\frac{(1-a)(1-b)}{(1+a)(1+b)}\,\beta_{p,q},
\]
where
\[
  a=2\lambda_{n+1}\sin^2\frac{p\pi}{2N},
  \qquad
  b=2\lambda_{n+1}\sin^2\frac{q\pi}{2N}.
\]
As in the Dirichlet case, for \(a,b\ge0\) not both zero, the modulus
is less than \(1\): convergence.

\begin{remark}
\note{Added beyond the notes.}
For \(p=q=0\) we get \(a=b=0\) and the eigenvalue \(1\); the
eigenvector \(\beta_{0,0}\) is the constant vector \(\mathbf 1\), the
null vector of \(A\).  So the iteration converges to a solution of
\(AU=2hG\); the additive constant in the limit is inherited from the
initial guess, as expected for the Neumann problem.
\end{remark}

%=====================================================================
\section{ADI Iteration for the Robin Problem}
%=====================================================================

\lecinfo{Lecture 40, 13 November 2026.}

Consider
\[
  \pdd u{x_1}+\pdd u{x_2}=0,\qquad 0\le x_1,x_2\le1,
\]
subject to
\begin{align*}
  \pd u{x_1}-p_1u&=f_0(x_2)\quad\text{on }x_1=0,
  &\pd u{x_1}+q_1u&=f_1(x_2)\quad\text{on }x_1=1,\\
  \pd u{x_2}-p_2u&=g_0(x_1)\quad\text{on }x_2=0,
  &\pd u{x_2}+q_2u&=g_1(x_1)\quad\text{on }x_2=1,
\end{align*}
with \(p_1,q_1,p_2,q_2>0\) and \(Nh=1\).  As for the Neumann problem,
all \((N+1)^2\) grid points are unknowns.  The difference system is
\[
  AU=hG,\qquad A=H_1+V_1,
\]
with \(H_1=\operatorname{diag}(E,\dots,E)\) block diagonal of order
\((N+1)^2\), and the Peaceman--Rachford iteration is
\begin{align*}
  \Bigl(\frac2{\lambda_{n+1}}I+H_1\Bigr)U^{n+1*}
  &=\Bigl(\frac2{\lambda_{n+1}}I-V_1\Bigr)U^n+hG,\\
  \Bigl(\frac2{\lambda_{n+1}}I+V_1\Bigr)U^{n+1}
  &=\Bigl(\frac2{\lambda_{n+1}}I-H_1\Bigr)U^{n+1*}+hG,
\end{align*}
where the steps are
\[
  \bigl(1-\tfrac12\lambda\delta_{x_1}^2\bigr)U_{l,m}^{n+1*}
  =\bigl(1+\tfrac12\lambda\delta_{x_2}^2\bigr)U_{l,m}^n,
  \qquad
  \bigl(1-\tfrac12\lambda\delta_{x_2}^2\bigr)U_{l,m}^{n+1}
  =\bigl(1+\tfrac12\lambda\delta_{x_1}^2\bigr)U_{l,m}^{n+1*}.
\]

\note{The matrices \(E\), \(V_1\), the vector \(G\) and the
convergence argument were left incomplete in lecture; they are
completed below.}

\subsection{\texorpdfstring{The matrices \(E\), \(V_1\) and the vector \(G\)}{The matrices E, V1 and the vector G}}

\step{Fictitious points}
At a point \((0,m)\) on \(x_1=0\), the central difference
\[
  \frac{u_{1,m}-u_{-1,m}}{2h}-p_1u_{0,m}=f_{0,m}
  \quad\Longrightarrow\quad
  u_{-1,m}=u_{1,m}-2hp_1u_{0,m}-2hf_{0,m}.
\]
So the \(x_1\)-part of the five-point formula (multiplied by \(-1\))
becomes
\[
  2u_{0,m}-u_{1,m}-u_{-1,m}=2(1+hp_1)\,u_{0,m}-2u_{1,m}+2hf_{0,m}.
\]
At a point \((N,m)\) on \(x_1=1\),
\[
  \frac{u_{N+1,m}-u_{N-1,m}}{2h}+q_1u_{N,m}=f_{1,m}
  \quad\Longrightarrow\quad
  u_{N+1,m}=u_{N-1,m}-2hq_1u_{N,m}+2hf_{1,m},
\]
\[
  2u_{N,m}-u_{N+1,m}-u_{N-1,m}=-2u_{N-1,m}+2(1+hq_1)\,u_{N,m}-2hf_{1,m}.
\]
The sides \(x_2=0\) and \(x_2=1\) are treated in the same way, with
\(p_2,q_2,g_0,g_1\).

\step{The matrices}
Ordering \(U\) with the \(x_1\) index running fastest, as for the
Neumann problem,
\begin{keybox}
\[
  E=\begin{bmatrix}
    2(1+hp_1)&-2&&&\\
    -1&2&-1&&\\
    &\ddots&\ddots&\ddots&\\
    &&-1&2&-1\\
    &&&-2&2(1+hq_1)
  \end{bmatrix}_{(N+1)\times(N+1)},
\]
\[
  V_1=\begin{bmatrix}
    2(1+hp_2)I&-2I&&&\\
    -I&2I&-I&&\\
    &\ddots&\ddots&\ddots&\\
    &&-I&2I&-I\\
    &&&-2I&2(1+hq_2)I
  \end{bmatrix}_{(N+1)^2\times(N+1)^2},
\]
\end{keybox}
where \(I\) is the identity of order \(N+1\).  The entry of \(G\) at
the grid point \((l,m)\) is
\[
  G_{l,m}=2\bigl(-f_{0,m}[l=0]+f_{1,m}[l=N]-g_{0,l}[m=0]+g_{1,l}[m=N]\bigr),
\]
where \([\,\cdot\,]\) is \(1\) if the condition holds and \(0\)
otherwise.  So \(G_{l,m}=0\) at interior points, and a corner receives
contributions from both of its sides.

\step{Limit of the iteration}
Exactly as for the Neumann problem, putting \(U^{n+1}=U^n=U\) gives
\[
  \frac4{\lambda_{n+1}}(H_1+V_1)U=\frac4{\lambda_{n+1}}hG
  \quad\Longrightarrow\quad
  (H_1+V_1)\,U=hG .
\]

\subsection{Convergence}

\(H_1=\operatorname{diag}(E,\dots,E)\) acts along \(x_1\) and \(V_1\)
along \(x_2\), so \(H_1V_1=V_1H_1\).  Let \(E\phi_p=\mu_p\phi_p\), and
let \(E_2\psi_q=\nu_q\psi_q\), where \(E_2\) is the
\((N+1)\times(N+1)\) analogue of \(E\) for the \(x_2\) direction
(with \(p_2,q_2\) in place of \(p_1,q_1\)).  Then the vectors
\(\beta_{p,q}\) with components
\((\beta_{p,q})_{l,m}=(\phi_p)_l\,(\psi_q)_m\) satisfy
\[
  H_1\beta_{p,q}=\mu_p\beta_{p,q},
  \qquad
  V_1\beta_{p,q}=\nu_q\beta_{p,q}.
\]

\step{The eigenvalues are positive}
Let \(W=\operatorname{diag}(\tfrac12,1,\dots,1,\tfrac12)\).  Then
\[
  WE=\begin{bmatrix}
    1+hp_1&-1&&\\
    -1&2&-1&\\
    &\ddots&\ddots&\ddots\\
    &&-1&1+hq_1
  \end{bmatrix}
\]
is symmetric, and for \(x\ne0\)
\[
  x^T(WE)x=hp_1x_0^2+hq_1x_N^2+\sum_{l=0}^{N-1}(x_{l+1}-x_l)^2>0,
\]
because if the sum vanishes then \(x\) is constant and
\(hp_1x_0^2>0\).  So \(WE\) is symmetric positive definite.  Since
\(E=W^{-1}(WE)\) is similar to
\(W^{-1/2}(WE)W^{-1/2}\), also symmetric positive definite, the
eigenvalues \(\mu_p\) of \(E\) are real and positive, and \(E\) has a
complete set of eigenvectors.  The same holds for \(\nu_q\).

\step{Spectral radius}
The iteration matrix
\[
  S_{n+1}=\Bigl(\frac2{\lambda_{n+1}}I+V_1\Bigr)^{-1}
  \Bigl(\frac2{\lambda_{n+1}}I+H_1\Bigr)^{-1}
  \Bigl(\frac2{\lambda_{n+1}}I-H_1\Bigr)\Bigl(\frac2{\lambda_{n+1}}I-V_1\Bigr)
\]
has the eigenvalues
\[
  \frac{\bigl(\frac2{\lambda_{n+1}}-\mu_p\bigr)\bigl(\frac2{\lambda_{n+1}}-\nu_q\bigr)}
       {\bigl(\frac2{\lambda_{n+1}}+\mu_p\bigr)\bigl(\frac2{\lambda_{n+1}}+\nu_q\bigr)}
  =\frac{(1-a)(1-b)}{(1+a)(1+b)},
  \qquad
  a=\frac{\lambda_{n+1}\mu_p}2>0,\quad b=\frac{\lambda_{n+1}\nu_q}2>0 .
\]
Each factor \(\bigl|\frac{1-a}{1+a}\bigr|<1\) for \(a>0\), so
\(\rho(S_{n+1})<1\) and the iteration converges for every
\(\lambda_{n+1}>0\).

\begin{remark}
Unlike the Neumann problem, all eigenvalues \(\mu_p+\nu_q\) of
\(A=H_1+V_1\) are positive, so \(A\) is nonsingular and the Robin
problem has a unique discrete solution.  There is no closed form for
\(\mu_p\) here; it depends on \(hp_1\) and \(hq_1\).
\end{remark}

%=====================================================================
\section{Hyperbolic Equations: The Wave Equation}
%=====================================================================

\lecinfo{Tablet notes, 16 November 2026.}

Consider the wave equation
\begin{equation}\label{eq:wave}
  \pdd ut=\pdd ux,\qquad 0<x<1,\ t>0,
\end{equation}
subject to
\[
  \text{IC: } u(x,0)=f(x),\quad \pd{u}{t}(x,0)=g(x);
  \qquad
  \text{BC: } u(0,t)=u(1,t)=0,\ t>0,
\]
on \([0,1]\times[0,\infty)\).

\subsection{The explicit scheme}

Use the difference approximations
\[
  \pdd ut=\frac1{k^2}\delta_t^2u+O(k^2),
  \qquad
  \pdd ux=\frac1{h^2}\delta_x^2u+O(h^2).
\]
Applying them to \eqref{eq:wave} at \((mh,nk)\),
\[
  \frac1{k^2}\delta_t^2U_m^n=\frac1{h^2}\delta_x^2U_m^n,
\]
that is,
\[
  \frac1{k^2}\bigl[U_m^{n+1}-2U_m^n+U_m^{n-1}\bigr]
  =\frac1{h^2}\bigl[U_{m+1}^n-2U_m^n+U_{m-1}^n\bigr].
\]

\begin{keybox}
\begin{definition}[Explicit scheme for the wave equation]
With \(r=k/h\),
\begin{equation}\label{eq:waveexp}
  U_m^{n+1}=2(1-r^2)\,U_m^n+r^2\bigl[U_{m+1}^n+U_{m-1}^n\bigr]-U_m^{n-1}.
\end{equation}
This is an explicit three-level formula.
\end{definition}
\end{keybox}

\begin{remark}
A three-level scheme is the minimum here, because the equation is
second order in \(t\).
\end{remark}

\step{Starting the computation}
To start the computation we need the data on two lines, \(t=0\) and
\(t=k\).  At \(t=0\), \(U_m^0=f_m\).  At \(t=k\), use a suitable
approximation to the initial condition on \(u_t\):
\[
  \pd{U_m^0}t=\frac1{2k}\bigl(U_m^1-U_m^{-1}\bigr)+O(k^2)=g_m
  \quad\Longrightarrow\quad
  U_m^{-1}=U_m^1-2k\,g_m .
\]
The extraneous values \(U_m^{-1}\) introduced by \eqref{eq:waveexp}
for \(n=0\) are removed using this relation.  Putting \(n=0\) in
\eqref{eq:waveexp},
\[
  U_m^1=2(1-r^2)f_m+r^2(f_{m+1}+f_{m-1})-U_m^1+2kg_m,
\]
\begin{keybox}
\[
  U_m^1=(1-r^2)f_m+\frac{r^2}2\bigl(f_{m+1}+f_{m-1}\bigr)+k\,g_m .
\]
\end{keybox}

\begin{exercise}
Show that the LTE of \eqref{eq:waveexp} is \(O(k^2+h^2)\).
\end{exercise}

\note{This was an exercise; solution not done in lecture.}

\begin{proof}[Solution]
For the exact solution,
\[
  \delta_t^2u=k^2u_{tt}+\frac{k^4}{12}u_{tttt}+O(k^6),
  \qquad
  \delta_x^2u=h^2u_{xx}+\frac{h^4}{12}u_{xxxx}+O(h^6).
\]
So
\[
  \frac{\delta_t^2u}{k^2}-\frac{\delta_x^2u}{h^2}
  =\bigl(u_{tt}-u_{xx}\bigr)+\frac1{12}\bigl(k^2u_{tttt}-h^2u_{xxxx}\bigr)+\cdots
  =\frac1{12}\bigl(k^2u_{tttt}-h^2u_{xxxx}\bigr)+\cdots,
\]
which is \(O(k^2+h^2)\).  Since \(u_{tttt}=u_{xxxx}\) for the wave
equation, the leading term is \(\frac{h^2}{12}(r^2-1)u_{xxxx}\); it
vanishes when \(r=1\).
\end{proof}

\begin{definition}[Consistency]
A difference approximation \eqref{eq:waveexp} for a hyperbolic
equation is consistent if
\[
  \frac{\text{truncation error}}{k^2}\to0\quad\text{as }h,k\to0,
\]
where the truncation error is that of the scheme written as in
\eqref{eq:waveexp} (i.e.\ multiplied through by \(k^2\)).
\end{definition}

\subsection{Stability: the CFL condition}

\begin{exercise}
Show that the von Neumann stability analysis gives the condition
\(r\le1\) for stability of \eqref{eq:waveexp}.
\end{exercise}

\note{This was an exercise; solution not done in lecture.}

\begin{proof}[Solution]
Put \(U_m^n=\xi^ne^{i\beta mh}\).  Since
\(\delta_x^2e^{i\beta mh}=-4\sin^2\frac{\beta h}2\,e^{i\beta mh}\),
the scheme gives, with \(s=\sin^2\frac{\beta h}2\in[0,1]\),
\[
  \xi-2+\xi^{-1}=-4r^2s
  \quad\Longrightarrow\quad
  \xi^2-2\bigl(1-2r^2s\bigr)\xi+1=0 .
\]
The product of the two roots is \(1\).  If the roots are real and
distinct, one of them has modulus \(>1\).  So we need complex (or
equal) roots, i.e.\
\[
  \bigl|1-2r^2s\bigr|\le1
  \iff 0\le r^2s\le1 .
\]
Then \(|\xi_1|=|\xi_2|=1\).  This must hold for all \(s\in[0,1]\),
so \(r^2\le1\), i.e.\ \(r\le1\).
\end{proof}

\begin{keybox}
The stability condition \(r=k/h\le1\) for the explicit formula
\eqref{eq:waveexp} is known as the \emph{CFL (Courant--Friedrichs--Lewy)
condition}.
\end{keybox}

\subsection{Implicit schemes}

If we replace \(U_m^n\) on the right-hand side of
\(\frac1{k^2}\delta_t^2U_m^n=\frac1{h^2}\delta_x^2U_m^n\) by the
weighted sum
\[
  \theta U_m^{n+1}+(1-2\theta)U_m^n+\theta U_m^{n-1},
\]
we get
\begin{keybox}
\begin{equation}\label{eq:waveimp}
  \delta_t^2U_m^n=r^2\delta_x^2\bigl[\theta U_m^{n+1}+(1-2\theta)U_m^n
  +\theta U_m^{n-1}\bigr],
  \qquad 0\le\theta\le1 .
\end{equation}
\end{keybox}
The difference scheme is of order \(O(k^2+h^2)\).

\step{Stability condition}
\begin{enumerate}[label=(\roman*)]
  \item \(\theta\ge\frac14\) gives unconditional stability.
  \item \(0\le\theta<\frac14\) gives conditional stability:
    \(0<r^2\le\dfrac1{1-4\theta}\).
\end{enumerate}

\note{Proof not done in lecture.}

\begin{proof}
With \(U_m^n=\xi^ne^{i\beta mh}\) and \(s=\sin^2\frac{\beta h}2\),
\[
  \xi-2+\xi^{-1}=-4r^2s\bigl[\theta\xi+(1-2\theta)+\theta\xi^{-1}\bigr],
\]
i.e.\
\[
  \bigl(1+4r^2s\theta\bigr)\xi^2-2\bigl(1-2r^2s(1-2\theta)\bigr)\xi
  +\bigl(1+4r^2s\theta\bigr)=0 .
\]
Again the product of the roots is \(1\), so stability requires
\(|1-2r^2s(1-2\theta)|\le1+4r^2s\theta\).  The inequality
\(1-2r^2s(1-2\theta)\le1+4r^2s\theta\) always holds.  The other one,
\(-1+2r^2s(1-2\theta)\le1+4r^2s\theta\), is
\[
  r^2s(1-4\theta)\le1 .
\]
If \(\theta\ge\frac14\) this holds for all \(r\).  If
\(\theta<\frac14\), it must hold for \(s=1\), giving
\(r^2\le1/(1-4\theta)\).  (For \(\theta=0\) this is the CFL condition
\(r\le1\).)
\end{proof}

\step{The case \(\theta=\frac14\)}
Since
\(\frac14U^{n+1}+\frac12U^n+\frac14U^{n-1}=U^n+\frac14\delta_t^2U^n\),
the value \(\theta=\frac14\) gives the von Neumann scheme
\[
  \bigl(1-\tfrac14r^2\delta_x^2\bigr)\delta_t^2U_m^n=r^2\delta_x^2U_m^n .
\]

\subsection{Two space dimensions}

Consider
\begin{equation}\label{eq:wave2}
  \pdd ut=\pdd u{x_1}+\pdd u{x_2},
\end{equation}
\[
  \text{IC: } u(x_1,x_2,0)=f(x_1,x_2),\quad
  \pd{u}{t}(x_1,x_2,0)=g(x_1,x_2),\qquad 0\le x_1,x_2\le1,
\]
\[
  \text{BC: } u=0\ \text{on the boundary of the unit square},\ t>0 .
\]

\step{An explicit scheme}
\[
  \delta_t^2U_{l,m}^n=r^2\bigl(\delta_{x_1}^2+\delta_{x_2}^2\bigr)U_{l,m}^n,
\]
or
\begin{keybox}
\[
  U_{l,m}^{n+1}=2(1-2r^2)\,U_{l,m}^n
  +r^2\bigl(U_{l-1,m}^n+U_{l,m-1}^n+U_{l,m+1}^n+U_{l+1,m}^n\bigr)
  -U_{l,m}^{n-1}.
\]
\end{keybox}
\begin{itemize}
  \item Order of the method: \(O(k^2+h^2)\).
  \item Stability: \(0<r\le\dfrac1{\sqrt2}\).
\end{itemize}

\note{Derivation of the stability condition not done in lecture.}
With \(U_{l,m}^n=\xi^ne^{i(\beta_1l+\beta_2m)h}\),
\(s_1=\sin^2\frac{\beta_1h}2\), \(s_2=\sin^2\frac{\beta_2h}2\),
\[
  \xi^2-2\bigl(1-2r^2(s_1+s_2)\bigr)\xi+1=0,
\]
and, as in one dimension, stability requires
\(r^2(s_1+s_2)\le1\) for all \(s_1,s_2\in[0,1]\), i.e.\
\(2r^2\le1\), \(r\le1/\sqrt2\).

\step{Implicit scheme}
Generalizing the implicit scheme \eqref{eq:waveimp} to
\eqref{eq:wave2},
\[
  \delta_t^2U_{l,m}^n=r^2\Bigl[\theta\bigl(\delta_{x_1}^2+\delta_{x_2}^2\bigr)U_{l,m}^{n+1}
  +(1-2\theta)\bigl(\delta_{x_1}^2+\delta_{x_2}^2\bigr)U_{l,m}^n
  +\theta\bigl(\delta_{x_1}^2+\delta_{x_2}^2\bigr)U_{l,m}^{n-1}\Bigr],
\]
or
\begin{keybox}
\[
  \Bigl[1-\theta r^2\bigl(\delta_{x_1}^2+\delta_{x_2}^2\bigr)\Bigr]\delta_t^2U_{l,m}^n
  =r^2\bigl(\delta_{x_1}^2+\delta_{x_2}^2\bigr)U_{l,m}^n .
\]
\end{keybox}
Stability:
\begin{enumerate}[label=(\roman*)]
  \item \(\theta\ge\frac14\): unconditionally stable.
  \item \(0\le\theta<\frac14\): conditionally stable,
    \(r^2\le\dfrac1{2(1-4\theta)}\).
\end{enumerate}
(The proof is the one-dimensional one with \(s\) replaced by
\(s_1+s_2\le2\).)  An ADI method can be developed for this scheme.

%=====================================================================
\section{The Heat Equation: Comparison of Methods}
%=====================================================================

\begin{example}\label{ex:heat}
Solve the heat conduction equation
\[
  \pd ut=\pdd ux,
  \qquad
  u(x,0)=\sin\pi x,\ 0\le x\le1,
  \qquad
  u(0,t)=u(1,t)=0,
\]
by
\begin{enumerate}[label=(\roman*)]
  \item the Crank--Nicolson scheme,
  \item the explicit method (Schmidt method),
  \item the Du Fort--Frankel method,
  \item the Douglas method.
\end{enumerate}
\end{example}

\subsection{The four schemes in one dimension}

With \(r=k/h^2\):
\begin{align*}
  \text{(i) Crank--Nicolson:}\quad
  &\bigl(1-\tfrac r2\delta_x^2\bigr)U_m^{n+1}=\bigl(1+\tfrac r2\delta_x^2\bigr)U_m^n,\\
  \text{(ii) Explicit (Schmidt):}\quad
  &U_m^{n+1}=\bigl(1+r\delta_x^2\bigr)U_m^n,\\
  \text{(iii) Du Fort--Frankel:}\quad
  &(1+2r)U_m^{n+1}=2r\bigl(U_{m+1}^n+U_{m-1}^n\bigr)+(1-2r)U_m^{n-1},\\
  \text{(iv) Douglas:}\quad
  &\Bigl[1-\tfrac12\bigl(r-\tfrac16\bigr)\delta_x^2\Bigr]U_m^{n+1}
   =\Bigl[1+\tfrac12\bigl(r+\tfrac16\bigr)\delta_x^2\Bigr]U_m^n .
\end{align*}

\begin{remark}
\note{Added beyond the notes.}
For comparison, the standard properties of these schemes are:

\begin{center}
\renewcommand{\arraystretch}{1.4}
\begin{tabular}{lll}
\toprule
Scheme & Truncation error & Stability\\
\midrule
Explicit (Schmidt) & \(O(k+h^2)\) & \(r\le\frac12\)\\
Crank--Nicolson & \(O(k^2+h^2)\) & unconditional\\
Du Fort--Frankel & \(O(k^2+h^2+k^2/h^2)\) & unconditional\\
Douglas & \(O(k^2+h^4)\) & unconditional\\
\bottomrule
\end{tabular}
\end{center}

Du Fort--Frankel is consistent only if \(k/h\to0\).  For the Douglas
scheme, with \(s=\sin^2\frac{\beta h}2\), the amplification factor is
\(\xi=\dfrac{1-2(r+\frac16)s}{1+2(r-\frac16)s}\); one checks
\(-1\le\xi\le1\) for all \(r>0\).
\end{remark}

\subsection{The corresponding schemes in two dimensions}

For \(u_t=u_{x_1x_1}+u_{x_2x_2}\):

\step{Crank--Nicolson (factored) and its split, Peaceman--Rachford ADI}
\[
  \bigl(1-\tfrac r2\delta_{x_1}^2\bigr)\bigl(1-\tfrac r2\delta_{x_2}^2\bigr)U_{l,m}^{n+1}
  =\bigl(1+\tfrac r2\delta_{x_1}^2\bigr)\bigl(1+\tfrac r2\delta_{x_2}^2\bigr)U_{l,m}^n,
\]
\begin{align*}
  \bigl(1-\tfrac r2\delta_{x_1}^2\bigr)U_{l,m}^{n+1*}
  &=\bigl(1+\tfrac r2\delta_{x_2}^2\bigr)U_{l,m}^n,\\
  \bigl(1-\tfrac r2\delta_{x_2}^2\bigr)U_{l,m}^{n+1}
  &=\bigl(1+\tfrac r2\delta_{x_1}^2\bigr)U_{l,m}^{n+1*}.
\end{align*}

\step{Douglas--Rachford ADI}
\begin{align*}
  \bigl(1-r\delta_{x_1}^2\bigr)U_{l,m}^{n+1*}&=\bigl(1+r\delta_{x_2}^2\bigr)U_{l,m}^n,\\
  \bigl(1-r\delta_{x_2}^2\bigr)U_{l,m}^{n+1}&=U_{l,m}^{n+1*}-r\delta_{x_2}^2U_{l,m}^n .
\end{align*}

\step{Douglas (high-order) scheme and its ADI split}
\[
  \Bigl(1-\tfrac12\bigl(r-\tfrac16\bigr)\delta_{x_1}^2\Bigr)
  \Bigl(1-\tfrac12\bigl(r-\tfrac16\bigr)\delta_{x_2}^2\Bigr)U_{l,m}^{n+1}
  =\Bigl(1+\tfrac12\bigl(r+\tfrac16\bigr)\delta_{x_1}^2\Bigr)
  \Bigl(1+\tfrac12\bigl(r+\tfrac16\bigr)\delta_{x_2}^2\Bigr)U_{l,m}^n,
\]
\begin{keybox}
\begin{align*}
  \Bigl(1-\tfrac12\bigl(r-\tfrac16\bigr)\delta_{x_1}^2\Bigr)U_{l,m}^{n+1*}
  &=\Bigl(1+\tfrac12\bigl(r+\tfrac16\bigr)\delta_{x_2}^2\Bigr)U_{l,m}^n,\\
  \Bigl(1-\tfrac12\bigl(r-\tfrac16\bigr)\delta_{x_2}^2\Bigr)U_{l,m}^{n+1}
  &=\Bigl(1+\tfrac12\bigl(r+\tfrac16\bigr)\delta_{x_1}^2\Bigr)U_{l,m}^{n+1*}.
\end{align*}
\end{keybox}
Eliminating \(U^{n+1*}\) (the operators commute) recovers the factored
Douglas scheme, exactly as for Peaceman--Rachford.  This split is
also known as the Mitchell--Fairweather scheme; in the notes it is
labelled ``Douglas--Rachford ADI''.

\subsection{Worked numerical illustration}

\note{This numerical example was added beyond the notes.}

Take \(h=\frac14\) and \(k=\frac1{64}\), so \(r=k/h^2=\frac14\).
The interior points are \(x=\frac14,\frac12,\frac34\), and
\[
  U^0=\bigl[\sin\tfrac\pi4,\ \sin\tfrac\pi2,\ \sin\tfrac{3\pi}4\bigr]
  =[0.707107,\ 1,\ 0.707107].
\]
The exact solution is \(u(x,t)=e^{-\pi^2t}\sin\pi x\).

\step{Key observation}
The grid vector \(\sin(\pi x_m)\) is an eigenvector of \(\delta_x^2\)
(with zero boundary values):
\[
  \delta_x^2\sin(\pi x_m)=-4\sin^2\frac{\pi h}2\,\sin(\pi x_m)
  =-4s\,\sin(\pi x_m),
  \qquad s=\sin^2\frac\pi8=0.146447 .
\]
Hence every two-level scheme gives \(U_m^n=\xi^n\sin(\pi x_m)\), with
\(\xi\) its amplification factor for this mode, and by symmetry
\(U_1^n=U_3^n\).  Here \(4rs=s=0.146447\).

\step{(ii) Explicit}
\[
  \xi=1-4rs=0.853553,
\]
\[
  U^1=[0.603553,\ 0.853553,\ 0.603553],\qquad
  U^2=[0.515165,\ 0.728553,\ 0.515165].
\]

\step{(i) Crank--Nicolson, solved directly}
With \(U_1=U_3=a\), \(U_2=b\), the system
\(-\frac r2U_{m-1}^{1}+(1+r)U_m^1-\frac r2U_{m+1}^1
=\frac r2U_{m-1}^0+(1-r)U_m^0+\frac r2U_{m+1}^0\) reads
\begin{align*}
  1.25a-0.125b&=0.75(0.707107)+0.125(1)=0.655330,\\
  -0.25a+1.25b&=0.75(1)+0.125(2\times0.707107)=0.926777 .
\end{align*}
The determinant is \(1.25^2-0.125\times0.25=1.53125\), so
\begin{align*}
  a&=\frac{1.25(0.655330)+0.125(0.926777)}{1.53125}=0.610619,\\
  b&=\frac{1.25(0.926777)+0.25(0.655330)}{1.53125}=0.863545 .
\end{align*}
This agrees with \(\xi=\frac{1-2rs}{1+2rs}=0.863545\).  A second step
gives \(U^2=[0.527297,\ 0.745710,\ 0.527297]\).

\step{(iv) Douglas}
Here \(\frac12(r-\frac16)=\frac1{24}\) and
\(\frac12(r+\frac16)=\frac5{24}\), so
\[
  \xi=\frac{1-4\cdot\frac5{24}s}{1+4\cdot\frac1{24}s}
  =\frac{0.877961}{1.024408}=0.857044,
\]
\[
  U^1=[0.606021,\ 0.857043,\ 0.606021],\qquad
  U^2=[0.519386,\ 0.734522,\ 0.519386].
\]

\step{(iii) Du Fort--Frankel}
It is a three-level scheme, so \(U^1\) must come from a two-level
scheme; take the explicit \(U^1\).  With \(r=\frac14\):
\(1.5\,U_m^2=0.5\bigl(U_{m+1}^1+U_{m-1}^1\bigr)+0.5\,U_m^0\).  At
\(x=\frac12\):
\[
  U_2^2=\frac{0.5(0.603553+0.603553)+0.5(1)}{1.5}=0.735702,
\]
and at \(x=\frac14\):
\(U_1^2=\frac{0.5(0+0.853553)+0.5(0.707107)}{1.5}=0.520220\).

\needspace{16\baselineskip}
\step{Comparison at \(x=\frac12\), \(t=2k=\frac1{32}\)}
The exact value is \(e^{-\pi^2/32}=0.734603\).

\begin{center}
\renewcommand{\arraystretch}{1.4}
\begin{tabular}{lcc}
\toprule
Method & \(U\bigl(\tfrac12,\tfrac1{32}\bigr)\) & Error\\
\midrule
Explicit & \(0.728553\) & \(6.0\times10^{-3}\)\\
Crank--Nicolson & \(0.745710\) & \(1.1\times10^{-2}\)\\
Du Fort--Frankel & \(0.735702\) & \(1.1\times10^{-3}\)\\
Douglas & \(0.734522\) & \(8.1\times10^{-5}\)\\
Exact & \(0.734603\) & --\\
\bottomrule
\end{tabular}
\end{center}

The Douglas scheme, of order \(O(k^2+h^4)\), is by far the most
accurate on this coarse mesh.  Crank--Nicolson is less accurate than
the explicit scheme here: both have an \(O(h^2)\) spatial error, and
at \(r=\frac14\) the explicit scheme's time error partly cancels its
space error, while for Crank--Nicolson the two errors add.

\end{document}
